\subsection{Core Artificial Intelligence Technologies}
\subsubsection{LLMs in Education}
LLMs, characterized by transformer architectures and self-attention mechanisms, have shifted educational technology from static content delivery toward context-aware generative systems capable of deep semantic understanding. In the aBridgeAI framework, LLMs serve as the primary engine for ``fading support'' scaffolding, dynamically adjusting the complexity of hints and explanations based on the learner's demonstrated competence. These models facilitate ``System-2 level engagement'' by generating automated questioning prompts that require students to explain their reasoning in steps, identify underlying assumptions, or solve problems through alternative methodologies. This approach prevents the ``AI as Autopilot'' failure mode, where a tool provides complete solutions that reduce germane load and hinder mastery; instead, the LLM enforces the productive mental work necessary for schema construction. Furthermore, LLMs enable the large-scale production of high-quality practice questions and quizzes, satisfying the need for a continuous supply of assessment materials in adaptive learning environments without the high cost and resource constraints of manual question construction \cite{kasneci2023chatgpt}.

Beyond generative scaffolding, LLMs fundamentally resolve the assessment bottlenecks prevalent in legacy LMS. Traditional platforms rely heavily on rigid RegEx or exact string matching, forcing a systemic dependence on MCQs that primarily stimulate recognition memory. LLMs, however, possess the high-dimensional semantic parsing capabilities required to reliably evaluate open-ended, production-based tasks. By analyzing the semantic intent of a student's constructed response---whether written or transcribed via voice---the LLM can accurately distinguish between genuine structural comprehension and superficial keyword mimicry, enabling nuanced, automated grading of short-answer questions and complex technical explanations.

Specifically within the aBridgeAI architecture, LLMs are operationalized to drive the context-aware mock interview module. This functionality requires the model to maintain conversational state and context over multi-turn dialogues, dynamically adapting its follow-up inquiries based on the learner's real-time articulation. To ensure the LLM strictly adheres to these pedagogical objectives rather than defaulting to a subservient, answer-providing assistant, the system employs advanced prompt engineering frameworks, including few-shot prompting and strict persona constraints. By functioning as an adaptive Socratic evaluator, the LLM continuously probes the boundaries of the student's knowledge, actively identifying and reporting the discrepancies between theoretical recall and practical application.

\subsubsection{Relational Reasoning: Overcoming Document Silos}
Traditional RAG pipelines often suffer from "document silos," where the system treats individual documents or text chunks as isolated units of information. This structural blindness limits relational reasoning, making it difficult for an AI to answer complex queries that require linking configuration rules, transactional dependencies, or conceptual hierarchies distributed across multiple sources\cite{lewis2020retrieval}. aBridgeAI addresses this through a structured lecture ingestion pipeline, where uploaded source materials — spanning documents, videos, and presentation files — are parsed and decomposed into typed knowledge units rather than stored as raw, undifferentiated text chunks. Each unit preserves its relational origin within the lesson it belongs to, allowing the system to reason across source boundaries rather than treating each file as a self-contained silo. Knowledge is represented not as a flat collection of passages but as a structured set of concepts, summaries, and examples that together form a coherent pedagogical context.

At the retrieval layer, the system combines parametric memory — the LLM's pre-trained knowledge — with non-parametric memory held in a dense vector index built over these structured materials, following the hybrid retrieval architecture described by Lewis et al. \cite{lewis2020retrieval}. This ensures that AI-generated assessments are grounded in verified, teacher-uploaded content rather than produced purely from the model's internal weights, reducing the risk of hallucination and improving factual specificity. The generation layer is further constrained by teacher-defined configuration. For quiz generation, instructors can supply additional contextual instructions that narrow the retrieval scope to lesson-relevant material. For interview generation, teacher-defined outcome requirements serve as semantic anchors that guide which regions of the knowledge index are consulted when producing sample questions — without requiring a pre-defined answer key from the teacher. This mirrors the principle of cognitive load management described by Sweller et al. (2019)\cite{Cog2019Sweller}, where scaffolded knowledge structures reduce extraneous cognitive load by narrowing the problem space to what is pedagogically relevant, benefiting both the learner and the generative process itself.
Together, these design choices position aBridgeAI's knowledge layer not as a collection of flat files but as a structured cognitive ecosystem — one where the relational context between materials, lessons, and assessment outcomes is preserved architecturally, enabling more coherent, traceable, and pedagogically grounded AI-generated content\cite{rag2025}.

RAG pipelines address the inherent limitations of purely parametric LLMs, such as their tendency to produce "hallucinations" and their inability to easily expand or revise stored world knowledge. By combining a pre-trained sequence-to-sequence model (parametric memory) with a dense vector index of verified instructional material (non-parametric memory), RAG anchors the AI’s output to specific, factual lesson contexts. This architecture ensures that the "supportive information" provided to a learner—which often involves high element interactivity and complex reasoning—is retrieved accurately from a ground-truth source before the LLM generates a response. The retrieval mechanism allows the system to be ``human-readable'' and ``human-writable'', enabling educators to update the knowledge base by simply editing the document index without retraining the underlying model parameters. Ultimately, this hybrid approach generates language that is more factual, specific, and diverse than parametric-only baselines, ensuring the learner is presented with information that is both reliable and contextually relevant to their current learning objective\cite{lewis2020retrieval}.

\subsubsection{Vector Databases and High-Dimensional Search}
Vector databases solve the problem of ``semantic isolation'' by embedding knowledge into high-dimensional vector spaces, where information is represented as mathematical coordinates based on semantic meaning. Unlike traditional keyword search, which relies on literal word overlap, high-dimensional search utilizes Maximum Inner Product Search (MIPS) and cosine similarity to identify conceptually related documents. This capability transforms the system from a collection of isolated data points into a ``dynamic, interdependent cognitive ecosystem'' where semantically related ``isolated elements'' are sequenced to manage the learner's working memory. When a student struggles with a concept, the vector database can retrieve ``just-in-time'' procedural information or supportive theoretical backgrounds by identifying the nearest neighbor concepts in the knowledge space. By mapping the dependencies between item difficulty, memory stability, and semantic relevance, high-dimensional search ensures that new information is integrated into a coherent mental model, effectively automating the ``isolated elements effect'' to prevent cognitive overload during the integration of complex causal systems\cite{lewis2020retrieval, Cog2019Sweller, mayer2003}.